\documentclass[11pt]{article}
\usepackage[margin=1in]{geometry}
\usepackage{amsmath,amssymb}
\usepackage{hyperref}
\usepackage{enumitem}
\usepackage{listings}

\begin{document}
\title{Implementation Notes for Superluminal Phenomena Modeling --- Pass 3}
\date{}
\maketitle

\section*{Goal}
Build a causally safe model where apparent superluminal signatures are observer-space transforms, while transport channels remain subluminal.

\section{Architecture}
\begin{enumerate}
\item Source preset: mass, spin, jet axis, distance, environment flags.
\item Event kernel: launch time, intrinsic speed, power, and channel energy split.
\item Observer transform: \(\beta_{\mathrm{app}}\), beaming, time-delay.
\item Multi-channel transport: EM, GW, CR, neutrino.
\item Output rendering: trajectories, light curves, posterior summaries.
\end{enumerate}

\section{Core equations}
\[
\beta_{\mathrm{app}}=\frac{\beta\sin\theta}{1-\beta\cos\theta},\qquad
t_{\mathrm{obs}}=t_{\mathrm{em}}+\frac{D-\mathbf{\hat n}\cdot\mathbf r_{\mathrm{em}}}{c}
\]
\[
F_{\nu,\mathrm{obs}}=\delta^{p}F'_{\nu'},\qquad
\delta=[\Gamma(1-\beta\cos\theta)]^{-1}
\]
\[
K_c(t)=A_c\mathbb{1}_{t\ge t_0}\left(\frac{t-t_0}{\tau_c}\right)e^{-(t-t_0)/\tau_c}
\]

\section{Python workflow}
\begin{itemize}
\item Core stack: \texttt{numpy}, \texttt{scipy}, \texttt{astropy}, \texttt{matplotlib}, \texttt{pandas}.
\item Bayesian stack: \texttt{numpyro}, \texttt{emcee}, \texttt{arviz}.
\item Suggested package modules:
\begin{itemize}
\item \texttt{src/core/geometry.py}
\item \texttt{src/transport/channels.py}
\item \texttt{src/observables/beaming.py}
\item \texttt{src/fit/objective.py}
\item \texttt{src/run/sweep.py}
\end{itemize}
\end{itemize}

\section{Julia workflow}
\begin{itemize}
\item \texttt{DifferentialEquations.jl}, \texttt{Turing.jl}, \texttt{DataFrames.jl}.
\item Use fast parametric sweeps and posterior sampling for posterior surfaces in \((\Gamma,\theta,\tau_c,p)\).
\end{itemize}

\section{MHD branch}
\begin{enumerate}
\item Run GRMHD/RMHD snapshot with HARM/PLUTO/BHAC.
\item Ray-project synthetic emissivity maps.
\item Feed map moments into apparent-motion and delay transforms.
\item Compare to MOJAVE-like trajectory and flux diagnostics.
\end{enumerate}

\section{Diagnostics}
\begin{itemize}
\item \(\beta_{\mathrm{app}}(t)\) tracks with causality envelope.
\item \(F_{\nu,\mathrm{obs}}(t)\) residuals vs multi-band data.
\item Corner plots over core parameters.
\item Channel-lag distributions and posterior predictive checks.
\end{itemize}

\section{Reproducibility protocol}
\begin{enumerate}
\item Version-pin all dependencies.
\item Archive random seeds and input data hashes.
\item Save config/run manifests with every experiment.
\item Keep `advanced_mode` (speculative couplings) disabled by default.
\end{enumerate}

\section{Suggested next step}
Create a single CLI entrypoint:
\texttt{simulate --source cenA --mode physical --sweep --backend python|julia}

\end{document}

